\medskip

Un paquet de lessive vide pèse 200~g. On y verse de la lessive.

On sait que 1~cm$^3$ de lessive pèse 1,5~g.

\medskip

\begin{enumerate}
\item Quelle est la masse totale d'un paquet de lessive (masse de la lessive et masse du paquet vide) contenant \np{1600}~ cm$^3$ de lessive ?
\item On considère la fonction $f$ qui à $x$ associe $1,5x + 200$.
	\begin{enumerate}
		\item Lorsque $x$ représente le volume de lessive en cm$^3$, que représente la valeur $f(x)$ ? 
		\item Représenter graphiquement la fonction $f$ dans un repère orthogonal.

On placera l'origine du repère en bas à gauche sur une feuille de papier millimétré. Sur l'axe des abscisses on prendra 1 cm pour 200~cm$^3$ et sur l'axe des ordonnées 1~cm pour $200$~g.
	\end{enumerate}
\item
	\begin{enumerate}
		\item En laissant les traits de construction apparents, trouver, par lecture graphique, le volume de lessive contenu dans un paquet de lessive de \np{2300}~g.
		\item Retrouver ce résultat par le calcul.
		\item Un paquet de lessive en forme de pavé de largeur 12~cm, de profondeur 8~cm et de
hauteur 15~cm peut-il contenir un tel volume ? 

\textbf{Argumenter la réponse en précisant la démarche.}
	\end{enumerate}
\end{enumerate}

\medskip
